\subsection{Questionnaire Instrument and Aggregate Responses}
\label{app:questionnaire_instrument}

The questionnaire was used only for task-design motivation, not for generating benchmark labels or experimental results. Tables~\ref{tab:questionnaire_choice_summary}--\ref{tab:questionnaire_expert_confirmation} summarize the translated aggregate report from the source spreadsheet, including all options, counts, percentages, and valid-response counts where available.

\begin{table*}[h]
\centering
\caption{Questionnaire aggregate responses for single-choice and multiple-choice items. Each option is shown as count (percentage) followed by the option text.}
\label{tab:questionnaire_choice_summary}
\scriptsize
\setlength{\tabcolsep}{3pt}
\begin{tabularx}{\textwidth}{p{0.05\linewidth}p{0.23\linewidth}X p{0.08\linewidth}}
\toprule
Item & Question & Options with count (percentage) & Valid N \\
\midrule
Q1 & What is your primary professional role? & \textbf{0} (0\%) University/research-institute researcher; \textbf{9} (52.94\%) University/research-institute student; \textbf{1} (5.88\%) Government/protected-area staff; \textbf{2} (11.76\%) Civil conservation organization staff/manager; \textbf{0} (0\%) Environmental consulting practitioner; \textbf{2} (11.76\%) Long-term volunteer/independent researcher; \textbf{3} (17.65\%) Other: [details] & 17 \\
Q2 & Do you use automated monitoring devices in your work, such as infrared cameras or acoustic recorders? & \textbf{4} (23.53\%) Frequent use; \textbf{8} (47.06\%) Occasional use; \textbf{5} (29.41\%) No current use; interested; \textbf{0} (0\%) No use; not familiar & 17 \\
Q4 & What are the biggest difficulties you currently face when processing monitoring data? & \textbf{7} (58.33\%) The data volume is too large to review manually; \textbf{10} (83.33\%) Poor photo quality makes identification difficult; \textbf{2} (16.67\%) There are too many species to identify and too many features to remember; \textbf{7} (58.33\%) Counting without double counting is difficult; \textbf{6} (50\%) Behavior understanding lacks a systematic method; \textbf{3} (25\%) Unusual cases are easy to miss (for example, hibernating animals active in winter); \textbf{3} (25\%) Linking monitoring data to environmental factors is difficult; \textbf{4} (33.33\%) Inconsistent formats are difficult to integrate; \textbf{0} (0\%) Other: [details] & 12 \\
Q6 & For the function you rated highest, please add detail 1: what specific problem do you most need this function to solve? & \textbf{6} (54.55\%) Save manual review time; \textbf{1} (9.09\%) Improve identification/counting accuracy; \textbf{1} (9.09\%) Find missed patterns or unusual cases; \textbf{3} (27.27\%) Standardize team workflow & 11 \\
Q7 & For the function you rated highest, please add detail 2: what form of output would you want this function to produce? & \textbf{0} (0\%) Direct final answer; \textbf{7} (63.64\%) Candidate list with confidence; \textbf{4} (36.36\%) Detailed analysis report; \textbf{0} (0\%) Mark suspicious/key points only & 11 \\
Q8 & For the function you rated highest, which type of error would be less acceptable to you? & \textbf{7} (63.64\%) False negative / missed real information; \textbf{4} (36.36\%) False positive / added review burden & 11 \\
Q9 & For the function you rated highest, how would you want its results to be handled? & \textbf{9} (81.82\%) Expert or user final confirmation; \textbf{2} (18.18\%) Auto-pass high confidence; review low confidence; \textbf{0} (0\%) Fully automatic & 11 \\
Q12 & What characteristics should a good intelligent foundation-model tool for ecological monitoring have? & \textbf{16} (94.12\%) High accuracy and few errors; \textbf{9} (52.94\%) Handles poor-quality photos; \textbf{9} (52.94\%) Explains judgments; \textbf{10} (58.82\%) Adapts to regions/species; \textbf{13} (76.47\%) Fits workflow and export formats; \textbf{13} (76.47\%) Supports correction and improvement; \textbf{1} (5.88\%) Other: [details] & 17 \\
\bottomrule
\end{tabularx}
\end{table*}

\begin{table}[h]
\centering
\caption{Ranked work activities from Q3.}
\label{tab:questionnaire_ranked_work}
\scriptsize
\setlength{\tabcolsep}{3pt}
\begin{tabularx}{\linewidth}{Xrrrr}
\toprule
Activity & Score & Rank 1 & Rank 2 & Rank 3 \\
\midrule
Species survey and inventory (identifying what species are present, where they occur, and how many there are) & 5.17 & 7(87.5\%) & 0(0\%) & 1(12.5\%) \\
Animal behavior research (what animals are doing and how they interact) & 3.67 & 3(50\%) & 2(33.33\%) & 1(16.67\%) \\
Science communication and public advocacy & 3.5 & 2(33.33\%) & 2(33.33\%) & 2(33.33\%) \\
Conservation project evaluation and management & 1.17 & 0(0\%) & 2(100\%) & 0(0\%) \\
Long-term ecological monitoring (tracking population changes and responses to climate change) & 1.08 & 0(0\%) & 1(50\%) & 1(50\%) \\
Other: & 1.08 & 0(0\%) & 1(50\%) & 1(50\%) \\
Addressing conservation threats (mitigating human-wildlife conflict, monitoring poaching, assessing engineering impacts) & 1 & 0(0\%) & 0(0\%) & 2(100\%) \\
Routine protected-area management (patrols, enforcement, habitat maintenance) & 0.58 & 0(0\%) & 1(100\%) & 0(0\%) \\
\bottomrule
\end{tabularx}
\end{table}

\begin{table*}[h]
\centering
\caption{Q5 function-value ratings. Top-box is the share of responses scoring 8--10; the distribution column reports percentages for scores 1--10.}
\label{tab:questionnaire_function_value}
\scriptsize
\setlength{\tabcolsep}{2.5pt}
\begin{tabularx}{\textwidth}{p{0.34\linewidth}p{0.08\linewidth}p{0.10\linewidth}X}
\toprule
Function & Mean & 8--10 & Score distribution (1--10) \\
\midrule
Animal counting & 8.75 & 66.66\% & 1:0\%, 2:0\%, 3:0\%, 4:0\%, 5:0\%, 6:0\%, 7:33.33\%, 8:8.33\%, 9:8.33\%, 10:50\% \\
Animal classification & 8.42 & 66.66\% & 1:0\%, 2:0\%, 3:8.33\%, 4:0\%, 5:0\%, 6:16.67\%, 7:8.33\%, 8:0\%, 9:8.33\%, 10:58.33\% \\
Fine-grained animal classification & 8.92 & 91.67\% & 1:0\%, 2:0\%, 3:0\%, 4:0\%, 5:0\%, 6:8.33\%, 7:0\%, 8:25\%, 9:25\%, 10:41.67\% \\
Behavior recognition & 8.33 & 75.00\% & 1:0\%, 2:0\%, 3:8.33\%, 4:0\%, 5:8.33\%, 6:0\%, 7:8.33\%, 8:16.67\%, 9:8.33\%, 10:50\% \\
Day/night and infrared activity-pattern analysis & 8 & 66.67\% & 1:0\%, 2:8.33\%, 3:0\%, 4:0\%, 5:8.33\%, 6:8.33\%, 7:8.33\%, 8:8.33\%, 9:16.67\%, 10:41.67\% \\
Seasonal plausibility assessment & 8.17 & 66.67\% & 1:0\%, 2:0\%, 3:0\%, 4:0\%, 5:0\%, 6:16.67\%, 7:16.67\%, 8:25\%, 9:16.67\%, 10:25\% \\
Unusual behavior discovery & 8.08 & 58.33\% & 1:0\%, 2:0\%, 3:0\%, 4:0\%, 5:0\%, 6:25\%, 7:16.67\%, 8:8.33\%, 9:25\%, 10:25\% \\
Natural occlusion grading & 8.25 & 75.00\% & 1:0\%, 2:0\%, 3:8.33\%, 4:0\%, 5:0\%, 6:8.33\%, 7:8.33\%, 8:25\%, 9:8.33\%, 10:41.67\% \\
Human-review routing (visibility triage) & 8.08 & 66.67\% & 1:0\%, 2:0\%, 3:0\%, 4:0\%, 5:16.67\%, 6:16.67\%, 7:0\%, 8:16.67\%, 9:8.33\%, 10:41.67\% \\
Low-light, white-flash, and IR imaging-mode difference tasks & 6.42 & 33.34\% & 1:0\%, 2:0\%, 3:8.33\%, 4:8.33\%, 5:25\%, 6:16.67\%, 7:8.33\%, 8:16.67\%, 9:0\%, 10:16.67\% \\
Temperature-conditioned activity analysis & 6.5 & 25.00\% & 1:0\%, 2:0\%, 3:8.33\%, 4:0\%, 5:33.33\%, 6:8.33\%, 7:25\%, 8:8.33\%, 9:0\%, 10:16.67\% \\
Region-specific ecological tasks & 6.25 & 33.33\% & 1:8.33\%, 2:0\%, 3:8.33\%, 4:8.33\%, 5:25\%, 6:0\%, 7:16.67\%, 8:8.33\%, 9:0\%, 10:25\% \\
Species interaction/co-occurrence analysis & 7.83 & 66.67\% & 1:0\%, 2:0\%, 3:0\%, 4:0\%, 5:25\%, 6:0\%, 7:8.33\%, 8:25\%, 9:16.67\%, 10:25\% \\
Long-term temporal pattern-shift mining & 7.33 & 58.33\% & 1:0\%, 2:0\%, 3:0\%, 4:8.33\%, 5:25\%, 6:8.33\%, 7:0\%, 8:25\%, 9:8.33\%, 10:25\% \\
Human disturbance/conflict-related tasks & 8.75 & 66.66\% & 1:0\%, 2:0\%, 3:0\%, 4:0\%, 5:0\%, 6:0\%, 7:33.33\%, 8:8.33\%, 9:8.33\%, 10:50\% \\
High-value rare event discovery & 8.17 & 66.67\% & 1:0\%, 2:0\%, 3:0\%, 4:0\%, 5:8.33\%, 6:8.33\%, 7:16.67\%, 8:16.67\%, 9:25\%, 10:25\% \\
\bottomrule
\end{tabularx}
\end{table*}

\begin{table*}[h]
\centering
\caption{Q10 preferred confirmation mode for each function. Percentages can sum above 100\% because this was a matrix multiple-choice item.}
\label{tab:questionnaire_expert_confirmation}
\scriptsize
\setlength{\tabcolsep}{3pt}
\begin{tabularx}{\textwidth}{Xrrr}
\toprule
Function & Expert required & Expert if low confidence & Fully automatic \\
\midrule
Animal counting & 0(0\%) & 8(66.67\%) & 6(50\%) \\
Animal classification & 8(66.67\%) & 4(33.33\%) & 2(16.67\%) \\
Fine-grained animal classification & 9(75\%) & 4(33.33\%) & 0(0\%) \\
Behavior recognition & 6(50\%) & 7(58.33\%) & 0(0\%) \\
Day/night and infrared activity-pattern analysis & 1(8.33\%) & 8(66.67\%) & 5(41.67\%) \\
Seasonal plausibility assessment & 2(16.67\%) & 10(83.33\%) & 1(8.33\%) \\
Unusual behavior discovery & 7(58.33\%) & 5(41.67\%) & 1(8.33\%) \\
Natural occlusion grading & 1(8.33\%) & 7(58.33\%) & 5(41.67\%) \\
Human-review routing (visibility triage) & 4(33.33\%) & 6(50\%) & 3(25\%) \\
Low-light, white-flash, and IR imaging-mode difference tasks & 1(8.33\%) & 6(50\%) & 6(50\%) \\
Temperature-conditioned activity analysis & 1(8.33\%) & 9(75\%) & 3(25\%) \\
Region-specific ecological tasks & 3(25\%) & 10(83.33\%) & 1(8.33\%) \\
Species interaction/co-occurrence analysis & 9(75\%) & 5(41.67\%) & 1(8.33\%) \\
Long-term temporal pattern-shift mining & 4(33.33\%) & 6(50\%) & 3(25\%) \\
Human disturbance/conflict-related tasks & 8(66.67\%) & 6(50\%) & 0(0\%) \\
High-value rare event discovery & 12(100\%) & 0(0\%) & 0(0\%) \\
\bottomrule
\end{tabularx}
\end{table*}

Q11 was a free-response item. The default aggregate spreadsheet states that detailed free-text answers are available on the detailed statistical-analysis page, but they are not included in the default report used for this paper draft.

